\chapter{Chronic Wounds}

\section*{Definition}
Wounds that fail to proceed through orderly healing stages within 4-6 weeks, including pressure ulcers, venous stasis ulcers, diabetic foot ulcers, and arterial insufficiency ulcers.

\section*{Pathophysiology}
Healing failure results from prolonged inflammation, impaired angiogenesis, bacterial biofilm formation, and matrix metalloproteinase imbalance. Underlying causes include ischemia (arterial disease), venous hypertension (venous insufficiency), neuropathy and pressure (diabetes), or prolonged pressure (immobility). Biofilm bacteria evade immune clearance and resist antibiotics.

\section*{Etiology}
\begin{itemize}
  \item Venous insufficiency (venous stasis ulcers)
  \item Peripheral arterial disease (ischemic ulcers)
  \item Diabetes mellitus (diabetic foot ulcers)
  \item Prolonged pressure (pressure ulcers/bedsores)
  \item Chronic infection or osteomyelitis
  \item Vasculitis or autoimmune disease
  \item Malignancy (basal cell carcinoma, squamous cell carcinoma)
  \item Radiation therapy damage
  \item Lymphedema
  \item Pyoderma gangrenosum
\end{itemize}

\section*{Indications for Evaluation}
Seek care if:
\begin{itemize}
  \item Severe or worsening symptoms
  \item Any non-healing wound persisting more than 4 weeks, wound with increasing pain, redness, swelling, or purulent drainage
  \item Accompanied by other concerning signs
\end{itemize}

\section*{Related Symptoms}
\begin{itemize}
  \item Leg swelling
  \item Skin discoloration
  \item Pain
  \item Fever
  \item Reduced mobility
\end{itemize}
\textbf{Note:} Always consider serious underlying conditions when evaluating persistent chronic wounds.

\section*{Self-Care / Home Management}
\begin{itemize}
  \item Rest and avoid activities that may aggravate the condition.
  \item Maintain adequate hydration and a balanced diet.
  \item Over-the-counter remedies may provide symptomatic relief (consult a pharmacist or doctor).
  \item Monitor symptoms and seek professional advice if they persist or worsen.
  \item Avoid self-medicating with prescription medications without medical guidance.
\end{itemize}

\section*{Diagnostic Approach}
\begin{itemize}
  \item A thorough history and physical examination are the first steps in evaluation.
  \item Laboratory tests (blood work, urinalysis) may be ordered based on clinical suspicion.
  \item Imaging studies (X-ray, ultrasound, CT, MRI) may be indicated when structural pathology is suspected.
  \item Specialist referral is arranged when the diagnosis remains uncertain or requires specific expertise.
\end{itemize}

\section*{Treatment / Management Overview}
\begin{itemize}
  \item Treatment targets the underlying cause identified through diagnostic evaluation.
  \item Supportive care and symptom management are provided while awaiting definitive therapy.
  \item Medications, lifestyle modifications, and physical therapies may be prescribed as appropriate.
  \item Follow-up ensures treatment effectiveness and allows timely adjustments to the care plan.
\end{itemize}

\clearpage